% ---------------------------------------------------------------------------
%  Shared preamble for the seven audience decks.
%
%  Each deck does, before the document body starts:
%     \newcommand{\capturedir}{../captures}
%     % ---------------------------------------------------------------------------
%  Shared preamble for the seven audience decks.
%
%  Each deck does, before the document body starts:
%     \newcommand{\capturedir}{../captures}
%     % ---------------------------------------------------------------------------
%  Shared preamble for the seven audience decks.
%
%  Each deck does, before the document body starts:
%     \newcommand{\capturedir}{../captures}
%     % ---------------------------------------------------------------------------
%  Shared preamble for the seven audience decks.
%
%  Each deck does, before the document body starts:
%     \newcommand{\capturedir}{../captures}
%     \input{persona-preamble}
%
%  \capturedir points at ../captures so both pools resolve:
%     \capture{24-scan.txt}              — the original v0.5.0 captures
%     \capture{persona/fe-01-scan-dist.txt} — captures recorded for these decks
% ---------------------------------------------------------------------------

\input{../common/evnx-preamble}

% ---- provenance ------------------------------------------------------------
% The original decks were pinned to tag v0.5.0 (85a52c3). Everything captured
% for the persona decks was re-run against 0.5.2, so the footline differs.
\renewcommand{\verifiedon}{\tiny\textcolor{evnxMuted}{verified against evnx 0.5.2 --- real run, not transcribed}}

% ---- audience badge on the title slide ------------------------------------
\newcommand{\audience}[1]{%
  \begin{center}\small\textcolor{evnxAccent}{\textbf{Audience: #1}}\end{center}}

% ---- verdict chips ---------------------------------------------------------
\newcommand{\solved}{\textcolor{evnxOk}{\ensuremath{\checkmark}}}
\newcommand{\partialmark}{\textcolor{evnxWarn}{\textbf{\textasciitilde}}}
\newcommand{\unsolved}{\textcolor{evnxErr}{\ensuremath{\times}}}

% ---- a correction block: something the source material got wrong -----------
\newenvironment{correction}[1]{%
  \begin{alertblock}{\small Correction --- #1}\small}{\end{alertblock}}

% ---- a trap block: works in the demo, fails in real use --------------------
\newenvironment{trap}[1]{%
  \setbeamercolor{block title}{bg=evnxWarn,fg=white}%
  \setbeamercolor{block body}{bg=evnxWarn!8}%
  \begin{block}{\small #1}\small}{\end{block}}

% ---- the persona card ------------------------------------------------------
\newcommand{\personacard}[4]{%
  \begin{block}{#1}
    \small
    \textbf{Stack:} #2 \\
    \textbf{Ships:} #3 \\
    \textbf{Loses time to:} #4
  \end{block}}

% ---- speaker note (rendered small and muted, not hidden) -------------------
\newcommand{\saythis}[1]{%
  \vspace{0.4em}{\footnotesize\textcolor{evnxMuted}{\textit{Say:} #1}}}

%
%  \capturedir points at ../captures so both pools resolve:
%     \capture{24-scan.txt}              — the original v0.5.0 captures
%     \capture{persona/fe-01-scan-dist.txt} — captures recorded for these decks
% ---------------------------------------------------------------------------

% ---------------------------------------------------------------------------
%  evnx presentation preamble — shared by all three decks
%
%  Engine: pdfLaTeX (also builds under XeLaTeX / LuaLaTeX).
%  Packages: beamer, listings, xcolor, amssymb, newunicodechar, booktabs,
%            fancyvrb, hyperref — all in texlive-latex-recommended + beamer.
%
%  \capturedir must be defined by the including document BEFORE \input of
%  this file, e.g.  \newcommand{\capturedir}{../captures}
% ---------------------------------------------------------------------------

\usetheme{Madrid}
\usecolortheme{seahorse}
\usefonttheme{professionalfonts}
\setbeamertemplate{navigation symbols}{}
\setbeamertemplate{caption}[numbered]
\setbeamerfont{frametitle}{size=\large}
\setbeamerfont{title}{size=\LARGE,series=\bfseries}

\usepackage[T1]{fontenc}
\usepackage[utf8]{inputenc}
\usepackage{lmodern}
\usepackage{amssymb}
\usepackage{booktabs}
\usepackage{xcolor}
\usepackage{listings}
\usepackage{fancyvrb}
\usepackage{newunicodechar}
\usepackage{tikz}
\usetikzlibrary{arrows.meta,positioning,shapes.geometric,fit,backgrounds}

% ---- palette ---------------------------------------------------------------
\definecolor{evnxInk}{HTML}{1F2933}
\definecolor{evnxAccent}{HTML}{2F6F4E}
\definecolor{evnxWarn}{HTML}{B7791F}
\definecolor{evnxErr}{HTML}{B03A2E}
\definecolor{evnxOk}{HTML}{2F6F4E}
\definecolor{evnxMuted}{HTML}{6B7280}
\definecolor{evnxTermBg}{HTML}{F5F5F2}
\definecolor{evnxRule}{HTML}{D1D5DB}

\setbeamercolor{structure}{fg=evnxAccent}
\setbeamercolor{alerted text}{fg=evnxErr}

% ---- glyph mapping for body text ------------------------------------------
% evnx emits exactly these 11 non-ASCII characters (enumerated from the
% captured logs, not guessed).
\newcommand{\evnxCheck}{\textcolor{evnxOk}{\ensuremath{\checkmark}}}
\newcommand{\evnxCross}{\textcolor{evnxErr}{\ensuremath{\times}}}
\newcommand{\evnxWarnSign}{\textcolor{evnxWarn}{\textbf{!}}}

\newunicodechar{✓}{\evnxCheck}
\newunicodechar{✗}{\evnxCross}
\newunicodechar{⚠}{\evnxWarnSign}
\newunicodechar{→}{\ensuremath{\rightarrow}}
\newunicodechar{↔}{\ensuremath{\leftrightarrow}}
\newunicodechar{─}{\textendash}
\newunicodechar{—}{---}
\newunicodechar{·}{\textperiodcentered}
\newunicodechar{…}{\ldots}
\newunicodechar{•}{\textbullet}
\newunicodechar{️}{}                % U+FE0F variation selector — invisible

% ---- terminal listing style ------------------------------------------------
% `literate` is required: newunicodechar does not apply inside listings.
\lstdefinestyle{evnxterm}{
  basicstyle=\ttfamily\scriptsize,
  backgroundcolor=\color{evnxTermBg},
  frame=single,
  rulecolor=\color{evnxRule},
  framesep=4pt,
  breaklines=true,
  breakatwhitespace=false,
  columns=fullflexible,
  keepspaces=true,
  showstringspaces=false,
  upquote=true,
  extendedchars=true,
  literate=
    {✓}{{\textcolor{evnxOk}{$\checkmark$}}}1
    {✗}{{\textcolor{evnxErr}{$\times$}}}1
    {⚠}{{\textcolor{evnxWarn}{\textbf{!}}}}1
    {→}{{$\rightarrow$}}1
    {↔}{{$\leftrightarrow$}}1
    {─}{{-}}1
    {—}{{---}}1
    {·}{{$\cdot$}}1
    {…}{{\ldots}}1
    {•}{{\textbullet}}1
    {️}{{}}0
    {é}{{\'e}}1
}

\lstdefinestyle{evnxcode}{
  style=evnxterm,
  basicstyle=\ttfamily\scriptsize,
  backgroundcolor=\color{white},
}

\lstset{style=evnxterm}

% ---- helper macros ---------------------------------------------------------
% \capture{file}        — include a real captured run verbatim
% \capturepart{f}{a}{b} — include lines a..b of a captured run
\newcommand{\capture}[1]{\lstinputlisting[style=evnxterm]{\capturedir/#1}}
\newcommand{\capturepart}[3]{%
  \lstinputlisting[style=evnxterm,firstline=#2,lastline=#3]{\capturedir/#1}}

% a labelled "real output" box
\newcommand{\realrun}[1]{%
  \begin{center}\scriptsize\textcolor{evnxMuted}{real run \textendash\ #1}\end{center}}

\newcommand{\cmd}[1]{\texttt{\textbf{#1}}}
\newcommand{\flag}[1]{\texttt{#1}}
\newcommand{\file}[1]{\texttt{#1}}

% exit-code chips
\newcommand{\exitzero}{\textcolor{evnxOk}{\texttt{0}}}
\newcommand{\exitone}{\textcolor{evnxWarn}{\texttt{1}}}
\newcommand{\exittwo}{\textcolor{evnxErr}{\texttt{2}}}

% a standard "verified against" footline for factual slides
\newcommand{\verifiedon}{\tiny\textcolor{evnxMuted}{verified against evnx 0.5.0 @ 9cfe75b}}

\usepackage[colorlinks=true,linkcolor=evnxAccent,urlcolor=evnxAccent]{hyperref}


% ---- provenance ------------------------------------------------------------
% The original decks were pinned to tag v0.5.0 (85a52c3). Everything captured
% for the persona decks was re-run against 0.5.2, so the footline differs.
\renewcommand{\verifiedon}{\tiny\textcolor{evnxMuted}{verified against evnx 0.5.2 --- real run, not transcribed}}

% ---- audience badge on the title slide ------------------------------------
\newcommand{\audience}[1]{%
  \begin{center}\small\textcolor{evnxAccent}{\textbf{Audience: #1}}\end{center}}

% ---- verdict chips ---------------------------------------------------------
\newcommand{\solved}{\textcolor{evnxOk}{\ensuremath{\checkmark}}}
\newcommand{\partialmark}{\textcolor{evnxWarn}{\textbf{\textasciitilde}}}
\newcommand{\unsolved}{\textcolor{evnxErr}{\ensuremath{\times}}}

% ---- a correction block: something the source material got wrong -----------
\newenvironment{correction}[1]{%
  \begin{alertblock}{\small Correction --- #1}\small}{\end{alertblock}}

% ---- a trap block: works in the demo, fails in real use --------------------
\newenvironment{trap}[1]{%
  \setbeamercolor{block title}{bg=evnxWarn,fg=white}%
  \setbeamercolor{block body}{bg=evnxWarn!8}%
  \begin{block}{\small #1}\small}{\end{block}}

% ---- the persona card ------------------------------------------------------
\newcommand{\personacard}[4]{%
  \begin{block}{#1}
    \small
    \textbf{Stack:} #2 \\
    \textbf{Ships:} #3 \\
    \textbf{Loses time to:} #4
  \end{block}}

% ---- speaker note (rendered small and muted, not hidden) -------------------
\newcommand{\saythis}[1]{%
  \vspace{0.4em}{\footnotesize\textcolor{evnxMuted}{\textit{Say:} #1}}}

%
%  \capturedir points at ../captures so both pools resolve:
%     \capture{24-scan.txt}              — the original v0.5.0 captures
%     \capture{persona/fe-01-scan-dist.txt} — captures recorded for these decks
% ---------------------------------------------------------------------------

% ---------------------------------------------------------------------------
%  evnx presentation preamble — shared by all three decks
%
%  Engine: pdfLaTeX (also builds under XeLaTeX / LuaLaTeX).
%  Packages: beamer, listings, xcolor, amssymb, newunicodechar, booktabs,
%            fancyvrb, hyperref — all in texlive-latex-recommended + beamer.
%
%  \capturedir must be defined by the including document BEFORE \input of
%  this file, e.g.  \newcommand{\capturedir}{../captures}
% ---------------------------------------------------------------------------

\usetheme{Madrid}
\usecolortheme{seahorse}
\usefonttheme{professionalfonts}
\setbeamertemplate{navigation symbols}{}
\setbeamertemplate{caption}[numbered]
\setbeamerfont{frametitle}{size=\large}
\setbeamerfont{title}{size=\LARGE,series=\bfseries}

\usepackage[T1]{fontenc}
\usepackage[utf8]{inputenc}
\usepackage{lmodern}
\usepackage{amssymb}
\usepackage{booktabs}
\usepackage{xcolor}
\usepackage{listings}
\usepackage{fancyvrb}
\usepackage{newunicodechar}
\usepackage{tikz}
\usetikzlibrary{arrows.meta,positioning,shapes.geometric,fit,backgrounds}

% ---- palette ---------------------------------------------------------------
\definecolor{evnxInk}{HTML}{1F2933}
\definecolor{evnxAccent}{HTML}{2F6F4E}
\definecolor{evnxWarn}{HTML}{B7791F}
\definecolor{evnxErr}{HTML}{B03A2E}
\definecolor{evnxOk}{HTML}{2F6F4E}
\definecolor{evnxMuted}{HTML}{6B7280}
\definecolor{evnxTermBg}{HTML}{F5F5F2}
\definecolor{evnxRule}{HTML}{D1D5DB}

\setbeamercolor{structure}{fg=evnxAccent}
\setbeamercolor{alerted text}{fg=evnxErr}

% ---- glyph mapping for body text ------------------------------------------
% evnx emits exactly these 11 non-ASCII characters (enumerated from the
% captured logs, not guessed).
\newcommand{\evnxCheck}{\textcolor{evnxOk}{\ensuremath{\checkmark}}}
\newcommand{\evnxCross}{\textcolor{evnxErr}{\ensuremath{\times}}}
\newcommand{\evnxWarnSign}{\textcolor{evnxWarn}{\textbf{!}}}

\newunicodechar{✓}{\evnxCheck}
\newunicodechar{✗}{\evnxCross}
\newunicodechar{⚠}{\evnxWarnSign}
\newunicodechar{→}{\ensuremath{\rightarrow}}
\newunicodechar{↔}{\ensuremath{\leftrightarrow}}
\newunicodechar{─}{\textendash}
\newunicodechar{—}{---}
\newunicodechar{·}{\textperiodcentered}
\newunicodechar{…}{\ldots}
\newunicodechar{•}{\textbullet}
\newunicodechar{️}{}                % U+FE0F variation selector — invisible

% ---- terminal listing style ------------------------------------------------
% `literate` is required: newunicodechar does not apply inside listings.
\lstdefinestyle{evnxterm}{
  basicstyle=\ttfamily\scriptsize,
  backgroundcolor=\color{evnxTermBg},
  frame=single,
  rulecolor=\color{evnxRule},
  framesep=4pt,
  breaklines=true,
  breakatwhitespace=false,
  columns=fullflexible,
  keepspaces=true,
  showstringspaces=false,
  upquote=true,
  extendedchars=true,
  literate=
    {✓}{{\textcolor{evnxOk}{$\checkmark$}}}1
    {✗}{{\textcolor{evnxErr}{$\times$}}}1
    {⚠}{{\textcolor{evnxWarn}{\textbf{!}}}}1
    {→}{{$\rightarrow$}}1
    {↔}{{$\leftrightarrow$}}1
    {─}{{-}}1
    {—}{{---}}1
    {·}{{$\cdot$}}1
    {…}{{\ldots}}1
    {•}{{\textbullet}}1
    {️}{{}}0
    {é}{{\'e}}1
}

\lstdefinestyle{evnxcode}{
  style=evnxterm,
  basicstyle=\ttfamily\scriptsize,
  backgroundcolor=\color{white},
}

\lstset{style=evnxterm}

% ---- helper macros ---------------------------------------------------------
% \capture{file}        — include a real captured run verbatim
% \capturepart{f}{a}{b} — include lines a..b of a captured run
\newcommand{\capture}[1]{\lstinputlisting[style=evnxterm]{\capturedir/#1}}
\newcommand{\capturepart}[3]{%
  \lstinputlisting[style=evnxterm,firstline=#2,lastline=#3]{\capturedir/#1}}

% a labelled "real output" box
\newcommand{\realrun}[1]{%
  \begin{center}\scriptsize\textcolor{evnxMuted}{real run \textendash\ #1}\end{center}}

\newcommand{\cmd}[1]{\texttt{\textbf{#1}}}
\newcommand{\flag}[1]{\texttt{#1}}
\newcommand{\file}[1]{\texttt{#1}}

% exit-code chips
\newcommand{\exitzero}{\textcolor{evnxOk}{\texttt{0}}}
\newcommand{\exitone}{\textcolor{evnxWarn}{\texttt{1}}}
\newcommand{\exittwo}{\textcolor{evnxErr}{\texttt{2}}}

% a standard "verified against" footline for factual slides
\newcommand{\verifiedon}{\tiny\textcolor{evnxMuted}{verified against evnx 0.5.0 @ 9cfe75b}}

\usepackage[colorlinks=true,linkcolor=evnxAccent,urlcolor=evnxAccent]{hyperref}


% ---- provenance ------------------------------------------------------------
% The original decks were pinned to tag v0.5.0 (85a52c3). Everything captured
% for the persona decks was re-run against 0.5.2, so the footline differs.
\renewcommand{\verifiedon}{\tiny\textcolor{evnxMuted}{verified against evnx 0.5.2 --- real run, not transcribed}}

% ---- audience badge on the title slide ------------------------------------
\newcommand{\audience}[1]{%
  \begin{center}\small\textcolor{evnxAccent}{\textbf{Audience: #1}}\end{center}}

% ---- verdict chips ---------------------------------------------------------
\newcommand{\solved}{\textcolor{evnxOk}{\ensuremath{\checkmark}}}
\newcommand{\partialmark}{\textcolor{evnxWarn}{\textbf{\textasciitilde}}}
\newcommand{\unsolved}{\textcolor{evnxErr}{\ensuremath{\times}}}

% ---- a correction block: something the source material got wrong -----------
\newenvironment{correction}[1]{%
  \begin{alertblock}{\small Correction --- #1}\small}{\end{alertblock}}

% ---- a trap block: works in the demo, fails in real use --------------------
\newenvironment{trap}[1]{%
  \setbeamercolor{block title}{bg=evnxWarn,fg=white}%
  \setbeamercolor{block body}{bg=evnxWarn!8}%
  \begin{block}{\small #1}\small}{\end{block}}

% ---- the persona card ------------------------------------------------------
\newcommand{\personacard}[4]{%
  \begin{block}{#1}
    \small
    \textbf{Stack:} #2 \\
    \textbf{Ships:} #3 \\
    \textbf{Loses time to:} #4
  \end{block}}

% ---- speaker note (rendered small and muted, not hidden) -------------------
\newcommand{\saythis}[1]{%
  \vspace{0.4em}{\footnotesize\textcolor{evnxMuted}{\textit{Say:} #1}}}

%
%  \capturedir points at ../captures so both pools resolve:
%     \capture{24-scan.txt}              — the original v0.5.0 captures
%     \capture{persona/fe-01-scan-dist.txt} — captures recorded for these decks
% ---------------------------------------------------------------------------

% ---------------------------------------------------------------------------
%  evnx presentation preamble — shared by all three decks
%
%  Engine: pdfLaTeX (also builds under XeLaTeX / LuaLaTeX).
%  Packages: beamer, listings, xcolor, amssymb, newunicodechar, booktabs,
%            fancyvrb, hyperref — all in texlive-latex-recommended + beamer.
%
%  \capturedir must be defined by the including document BEFORE \input of
%  this file, e.g.  \newcommand{\capturedir}{../captures}
% ---------------------------------------------------------------------------

\usetheme{Madrid}
\usecolortheme{seahorse}
\usefonttheme{professionalfonts}
\setbeamertemplate{navigation symbols}{}
\setbeamertemplate{caption}[numbered]
\setbeamerfont{frametitle}{size=\large}
\setbeamerfont{title}{size=\LARGE,series=\bfseries}

\usepackage[T1]{fontenc}
\usepackage[utf8]{inputenc}
\usepackage{lmodern}
\usepackage{amssymb}
\usepackage{booktabs}
\usepackage{xcolor}
\usepackage{listings}
\usepackage{fancyvrb}
\usepackage{newunicodechar}
\usepackage{tikz}
\usetikzlibrary{arrows.meta,positioning,shapes.geometric,fit,backgrounds}

% ---- palette ---------------------------------------------------------------
\definecolor{evnxInk}{HTML}{1F2933}
\definecolor{evnxAccent}{HTML}{2F6F4E}
\definecolor{evnxWarn}{HTML}{B7791F}
\definecolor{evnxErr}{HTML}{B03A2E}
\definecolor{evnxOk}{HTML}{2F6F4E}
\definecolor{evnxMuted}{HTML}{6B7280}
\definecolor{evnxTermBg}{HTML}{F5F5F2}
\definecolor{evnxRule}{HTML}{D1D5DB}

\setbeamercolor{structure}{fg=evnxAccent}
\setbeamercolor{alerted text}{fg=evnxErr}

% ---- glyph mapping for body text ------------------------------------------
% evnx emits exactly these 11 non-ASCII characters (enumerated from the
% captured logs, not guessed).
\newcommand{\evnxCheck}{\textcolor{evnxOk}{\ensuremath{\checkmark}}}
\newcommand{\evnxCross}{\textcolor{evnxErr}{\ensuremath{\times}}}
\newcommand{\evnxWarnSign}{\textcolor{evnxWarn}{\textbf{!}}}

\newunicodechar{✓}{\evnxCheck}
\newunicodechar{✗}{\evnxCross}
\newunicodechar{⚠}{\evnxWarnSign}
\newunicodechar{→}{\ensuremath{\rightarrow}}
\newunicodechar{↔}{\ensuremath{\leftrightarrow}}
\newunicodechar{─}{\textendash}
\newunicodechar{—}{---}
\newunicodechar{·}{\textperiodcentered}
\newunicodechar{…}{\ldots}
\newunicodechar{•}{\textbullet}
\newunicodechar{️}{}                % U+FE0F variation selector — invisible

% ---- terminal listing style ------------------------------------------------
% `literate` is required: newunicodechar does not apply inside listings.
\lstdefinestyle{evnxterm}{
  basicstyle=\ttfamily\scriptsize,
  backgroundcolor=\color{evnxTermBg},
  frame=single,
  rulecolor=\color{evnxRule},
  framesep=4pt,
  breaklines=true,
  breakatwhitespace=false,
  columns=fullflexible,
  keepspaces=true,
  showstringspaces=false,
  upquote=true,
  extendedchars=true,
  literate=
    {✓}{{\textcolor{evnxOk}{$\checkmark$}}}1
    {✗}{{\textcolor{evnxErr}{$\times$}}}1
    {⚠}{{\textcolor{evnxWarn}{\textbf{!}}}}1
    {→}{{$\rightarrow$}}1
    {↔}{{$\leftrightarrow$}}1
    {─}{{-}}1
    {—}{{---}}1
    {·}{{$\cdot$}}1
    {…}{{\ldots}}1
    {•}{{\textbullet}}1
    {️}{{}}0
    {é}{{\'e}}1
}

\lstdefinestyle{evnxcode}{
  style=evnxterm,
  basicstyle=\ttfamily\scriptsize,
  backgroundcolor=\color{white},
}

\lstset{style=evnxterm}

% ---- helper macros ---------------------------------------------------------
% \capture{file}        — include a real captured run verbatim
% \capturepart{f}{a}{b} — include lines a..b of a captured run
\newcommand{\capture}[1]{\lstinputlisting[style=evnxterm]{\capturedir/#1}}
\newcommand{\capturepart}[3]{%
  \lstinputlisting[style=evnxterm,firstline=#2,lastline=#3]{\capturedir/#1}}

% a labelled "real output" box
\newcommand{\realrun}[1]{%
  \begin{center}\scriptsize\textcolor{evnxMuted}{real run \textendash\ #1}\end{center}}

\newcommand{\cmd}[1]{\texttt{\textbf{#1}}}
\newcommand{\flag}[1]{\texttt{#1}}
\newcommand{\file}[1]{\texttt{#1}}

% exit-code chips
\newcommand{\exitzero}{\textcolor{evnxOk}{\texttt{0}}}
\newcommand{\exitone}{\textcolor{evnxWarn}{\texttt{1}}}
\newcommand{\exittwo}{\textcolor{evnxErr}{\texttt{2}}}

% a standard "verified against" footline for factual slides
\newcommand{\verifiedon}{\tiny\textcolor{evnxMuted}{verified against evnx 0.5.0 @ 9cfe75b}}

\usepackage[colorlinks=true,linkcolor=evnxAccent,urlcolor=evnxAccent]{hyperref}


% ---- provenance ------------------------------------------------------------
% The original decks were pinned to tag v0.5.0 (85a52c3). Everything captured
% for the persona decks was re-run against 0.5.2, so the footline differs.
\renewcommand{\verifiedon}{\tiny\textcolor{evnxMuted}{verified against evnx 0.5.2 --- real run, not transcribed}}

% ---- audience badge on the title slide ------------------------------------
\newcommand{\audience}[1]{%
  \begin{center}\small\textcolor{evnxAccent}{\textbf{Audience: #1}}\end{center}}

% ---- verdict chips ---------------------------------------------------------
\newcommand{\solved}{\textcolor{evnxOk}{\ensuremath{\checkmark}}}
\newcommand{\partialmark}{\textcolor{evnxWarn}{\textbf{\textasciitilde}}}
\newcommand{\unsolved}{\textcolor{evnxErr}{\ensuremath{\times}}}

% ---- a correction block: something the source material got wrong -----------
\newenvironment{correction}[1]{%
  \begin{alertblock}{\small Correction --- #1}\small}{\end{alertblock}}

% ---- a trap block: works in the demo, fails in real use --------------------
\newenvironment{trap}[1]{%
  \setbeamercolor{block title}{bg=evnxWarn,fg=white}%
  \setbeamercolor{block body}{bg=evnxWarn!8}%
  \begin{block}{\small #1}\small}{\end{block}}

% ---- the persona card ------------------------------------------------------
\newcommand{\personacard}[4]{%
  \begin{block}{#1}
    \small
    \textbf{Stack:} #2 \\
    \textbf{Ships:} #3 \\
    \textbf{Loses time to:} #4
  \end{block}}

% ---- speaker note (rendered small and muted, not hidden) -------------------
\newcommand{\saythis}[1]{%
  \vspace{0.4em}{\footnotesize\textcolor{evnxMuted}{\textit{Say:} #1}}}
